\section{Prospective image-to-text phase intervention}
\label{sec:crossmodal-phase-candidate}

We define a candidate that changes the conditional distribution over text keys
while preserving the deployed image-key probabilities and total text probability
at each image-query row, in real arithmetic. The hypothesis is that cross-modal
routing may benefit from a different rotary geometry than image-to-image
attention. The native capability screen supplies no primary semantic failure
that this candidate is known to repair.

For an image-token grid of height $H$ and width $W$, use the fixed reference map
\begin{equation}
 \Phi_{H,W}(p_0,p_h,p_w)
 =\left(p_0,\frac{63p_h}{H-1},\frac{63p_w}{W-1}\right).
 \label{eq:crossmodal-phase-map}
\end{equation}
Thus a $256\times256$ grid maps its spatial endpoints to $0$ and $63$;
the map is the identity at $64\times64$. The reference denotes our nominal
$1024^2$ setting, not a verified boundary of the model's training coordinates.
At each block we rotate the same normalized pre-RoPE image query using base
frequencies with $\theta=10000$ and axis dimensions $(16,56,56)$ at
$\Phi_{H,W}(p)$. We retain the actual pair-equal amplitude gains, text keys and
all values. Deployed SEGA geometry~\citep{sega2026} still determines image-key
logits. The intervention covers all 24 heads, all 57 joint and single blocks,
all 28 denoising steps and all 512 text slots, including padding. Runtime text
IDs must be zero. Text-query outputs remain unchanged within that block;
later text activations can change through earlier modified image states.

For one image query, let $a_j$ and $b_k$ be the deployed, scaled logits for text
keys $T$ and image keys $I$. Define $Z_T=\sum_{j\in T}e^{a_j}$,
$Z_I=\sum_{k\in I}e^{b_k}$, and $m=Z_T/(Z_T+Z_I)$. If $a'_j$ are the
transported-phase text logits, set
\begin{equation}
 \widetilde a_j=a'_j+\operatorname{LSE}(a)-\operatorname{LSE}(a'),
 \qquad b'_k=b_k.
 \label{eq:crossmodal-fixed-mass}
\end{equation}
The text partition remains $Z_T$, so every image-key probability and total text
mass $m$ are unchanged. With $u_T=\operatorname{softmax}(a)V_T$ and
$u'_T=\operatorname{softmax}(a')V_T$, the equivalent output correction is
\begin{equation}
 y'=y+m(u'_T-u_T).
 \label{eq:crossmodal-output-correction}
\end{equation}
These are elementary softmax identities. They impose a local constraint and
do not establish an improvement in the sampling trajectory.

The unnormalized control uses the same phase map but permits text mass to
change. It computes the original image-only conditional output $u_I$ and
log-partition $\ell_I$, and the new text-only output $u'_T$ and
log-partition $\ell'_T$. For $\ell=\operatorname{logaddexp}(\ell_I,\ell'_T)$,
\begin{equation}
 y_{\mathrm{unc}}=e^{\ell_I-\ell}u_I+e^{\ell'_T-\ell}u'_T.
\end{equation}
The original joint call supplies unchanged text-query outputs. This stable
mixture avoids recovering a small image contribution by subtracting rounded
text output from a BF16 joint output. Its extra image-only pass is a declared
cost. At 65,536 image tokens and 512 text slots, the fixed-mass arm adds
approximately $1.538\%$ of baseline joint-attention QK/AV arithmetic through
two text branches; the unnormalized control adds approximately $99.225\%$
through image-only and new-text branches. These counts exclude rotations,
buffers and kernel overhead, and are not latency predictions.

All three arms use a pinned private cuDNN output-and-log-sum-exp backend with
no fallback. The protocol requires tiny BF16 kernel comparisons against FP64
oracles and a production-shaped text-branch support probe before model loading.
Every full-resolution call must satisfy the backend's output/LSE contract.
A four-step $1024^2$ smoke compares the deployed attention path with the new
baseline, accepting endpoint relative $L_2$ error at most $10^{-3}$ and recording
cosine diagnostically. Both native phase arms must reproduce the matched
baseline endpoint bytes. Native execution shares the baseline path; separate
non-bypassed CPU oracles check the phase and partition formulas. The fixed-mass
correction accumulates in FP32 before its BF16 output cast. Real-arithmetic
conservation therefore remains distinct from floating-point equivalence.

Positional interpolation~\citep{chen2023pi}, cross-modal geometry separation
in Circle-RoPE,\footnote{Wang et al., \emph{Circle-RoPE: Cone-like Decoupled
Rotary Positional Embedding for Vision-Language Models},
\href{https://arxiv.org/abs/2505.16416v3}{arXiv:2505.16416v3}.}
query-dependent phase scales in Phase-Aligned RoPE,\footnote{Wu et al.,
\emph{Phase-Aligned RoPE for Mixed-Resolution Diffusion Transformer},
\href{https://arxiv.org/abs/2511.19778v2}{arXiv:2511.19778v2}.}
and token-group rotary edits in Untwisting RoPE\footnote{Mikaeili et al.,
\emph{Untwisting RoPE: Frequency Control for Shared Attention in DiTs},
\href{https://arxiv.org/abs/2602.05013v1}{arXiv:2602.05013v1}.}
are established ingredients. TACA already changes relative text-token odds
through modality-specific temperature,\footnote{Lv et al., \emph{Rethinking
Cross-Modal Interaction in Multimodal Diffusion Transformers},
\href{https://arxiv.org/abs/2506.07986v3}{arXiv:2506.07986v3}.}
and TIDE studies text/image imbalance~\citep{tide2026}. The proposed composition
isolates phase changes while holding text mass fixed; that constraint also
prevents it from correcting pure text-mass dilution. Neither novelty nor a
quality advantage has been established.

The frozen development pilot specifies exactly six $4096^2$ images: two existing
captions, seed 0, 28 steps and guidance 4.5 under the matched baseline,
fixed-mass and unnormalized arms. All use the same Krea backbone in BF16 and
must share exact actual initial transformer inputs within each prompt and
resolution. Both captions already satisfy their primary predicates, so this
pilot tests regression, detail and cost with little primary-semantic headroom.
Its scope excludes a seed search and fresh confirmation data; any readout must
retain all six images and report the unnormalized control's extra cost.
